\section{What the trajectories teach}
The most defensible shared lesson is that an economic action changes available options, and that statements about those changes can be compared with subsequent conduct. Finance provides a particularly clear correspondence. The state specifies a debt and resources, an initial statement can include collateral not yet listed, and the applied sequence either reaches payment or stops. In bargaining the same idea enters the language, but the goal of winning does not assign an observable value to every proposed exchange. This difference in endpoint clarity explains the different evidential weight we give the two results.

The finance evidence is positive even though the environment supports it with refreshed menus. Agents repeatedly continue liquidation when the immediately visible mortgage list is insufficient, and the entry statements usually name what selling will release. That is a concrete accomplishment under the evaluated interface. The corpus also contains correct local aggregate budgeting and sustained execution, including the \$14 sequence. Stating the assistance does not erase the behavior; it specifies the task accomplished. A claim that the same sequence demonstrates deep search would require evidence about dependency depth or performance under reduced feedback.

The two premature declarations clarify which decision is consequential. An agent that concludes its assets are insufficient can voluntarily stop before discovering the expanded menu. Successful entry statements instead preserve the possibility of that later action. This observation motivates measuring resource inclusion separately from action validity and episode completion. It also suggests why counting calls is a poor planning metric: many calls can follow one anticipated dependency, while one declaration can embody an incorrect global resource assessment. The retained records measure that contrast without requiring claims about a hidden search algorithm.

Bargaining demonstrates different degrees of economic assessment. The G15 response identifies rival building cash and follows through by rejecting three versions of a mutually completing trade. The G2 pricing rationale includes a correct cash remainder and an incorrect house count. G19 Sol corrects a false premise while accepting favorable terms. Thus expressed concern, factual accuracy and economic choice do not collapse into one variable. A benchmark that scores any rival-oriented sentence as strategic success would combine distinct achievements and overlook useful correction.

Cash is also relational. A buyer's payment can finance the seller's construction; the same payment can preserve a blocker against a third party; withholding it can accelerate the final opponent's elimination. The G7 construction continuation, G13 estate pitch and G18 survival denial instantiate those different stated stakes. This helps explain why a universal interpretation of cash-poor behavior is inadequate. The strategic relevance depends on who receives cash, what it enables, and who gains if the recipient fails. These are observable economic relations even when their optimal weight in a decision is unknown.

Communication adds information about those relations but does not establish persuasion. A public ceiling below a private contingency is a documented bluff under the task's permission. A selective financing claim can reach a buyer and precede acceptance without identifying the effect of the omission. A public call to withhold construction finance can be followed by the speaker's refusal while another player pursues the opposite trade. Keeping these different outcomes visible makes the incidents useful for studying negotiation without turning them into a catalogue of successful manipulation.

The complete offer table likewise changes what can be learned. Counters are frequent, and repricing a noncash bundle occurs for both completing and non-completing requests. A rejection-only summary omits much of how agents negotiate. The proposer strata then show why a pooled response difference is partly a property of the proposals encountered. Recipient concern may be real while the aggregate comparison remains confounded by compensation and counterpart behavior. The separate 61/74 population confirms expressed concern elsewhere, rather than resolving that identification problem.

These distinctions identify useful future targets without supplying their scores. A prospective comparison could vary rival development opportunities while documenting unavoidable changes to the recipient's own terms. A local financing comparison could measure payment and residual assets under the same menus and specified action ordering. Changing a word such as monopoly while holding the state fixed would answer a different question from changing the holdings that make development possible. None of those outcomes is inferred from the present cohort, and retrospective stability does not replace their prospective definition.

\section{Related work}
Neural negotiation research has long separated language generation from agreement and payoff. Lewis et al.~\cite{lewis2017deal} use multi-issue bargaining with hidden reward functions and evaluate dialogue rollouts. NegotiationArena studies multi-turn resource allocation, trading and price negotiation, including changes in tactics~\cite{bianchi2024negotiationarena}. Those controlled endpoints differ from a complete game's persistent and uncertain consequences. Our observations add a connection between exchanges and later finance, without providing a payoff reference of comparable strength for each bargain.

Guhe and Lascarides~\cite{guhe2014persuasion} already study trade embargoes and persuasion within a larger resource game. Their controlled symbolic-agent simulations examine when successful persuasion changes winning chances. We do not claim novelty for strategic withholding, nor infer its effectiveness from a refusal. Economic externality theory also formalizes why an asset's buyer can affect non-buyers~\cite{jehiel1996externalities}. G13 supplies an identity-dependent stated concern, not verification of an equilibrium theorem. One implication worth testing is whether willingness to finance a blocker changes with the identity of the creditor that would benefit from its failure.

Persistent LLM economic behavior is studied in Vending-Bench and state-linked communication research~\cite{backlund2025vendingbench,li2026arena}. AgentMonopoly examines model and persona variation, while FinBoardBench studies financial management across board games~\cite{wu2026agentmonopoly,hu2026finboardbench}. Playing Monopoly, counteroffering and managing liquidity are therefore not novel contributions by themselves. The present contribution is the complete response and debt accounting, joined to prospective resource statements and economic continuations in one cohort. Cross-study win rates would compare different rules, opponents and interfaces.

Counterparty-modeling research distinguishes information about another negotiator from using it effectively~\cite{cosentino2026counterparty}. Explanation-faithfulness work supplies a parallel warning about treating returned rationales as computation~\cite{turpin2023unfaithful,lanham2023faithfulness}. Our measurements respect both distinctions. They establish the content and timing of statements, some factual correctness, and executed local goals. They do not turn a plausible rationale into a general measure of strategy quality.

\section{Limitations}
The cohort contains one interacting roster, five seeds and four cyclic rotations, with no repeated seed/rotation cell. Relative opponent order stays fixed. Wins are 8/7/3/2 for Sol/Gemini/Opus/Grok and 7/7/3/3 by seat. Sol wins one entire block from every seat; another block is won from seat two by each system. These observations expose the difficulty of separating player ability, chance and opponent behavior. They do not quantify their causal shares.

All contrasts and case selections are retrospective. The sixteen-episode subset is defined by a resource inequality, but its prominence and semantic reading are also retrospective. Neither repeated arithmetic nor resampling validates its interpretation on new data. Five blocks poorly characterize tail behavior, provider changes or alternative rosters. Removing a system's rows changes a descriptive population without removing its historical influence on everyone else's state.

Unequal recovery paths, content/call disagreements and faulty historical summaries materially limit model-level interpretation. The prompts invite deception and provide incomplete board-value information. Semantic judgments are LLM-assisted, with no human annotation panel; related readers are not independent human validation. The necessary-case texts are first-response observations with different channels and varying numerical precision. There is no executed baseline, controlled communication intervention, trade-value reference, memory ablation or held-out cohort. Consequently, no conclusion establishes optimal strategy, causal persuasion, durable learning or general model rankings.

\section{Conclusion}
The games document agents assessing changing economic opportunities and sometimes carrying those assessments into successful action. Fourteen of sixteen obligations requiring newly released collateral reach payment, usually after that collateral appears in the entry statement. Bargaining documents rival-oriented concerns and a form-dependent rejection pattern whose magnitude depends on proposer composition. Contextual incidents show why those statements, factual checks and subsequent choices must remain separate. Following the whole trajectory makes local financial execution verifiable while preserving the uncertainty surrounding negotiation value and general strategic ability.
